\section{Preserving the later boundary conditions}
\label{sec:geometry-preservation}

We complete Proposition~\ref{prop:primitive-addition}.  The analytic
construction in Proposition~\ref{prop:primitive-realization} supplies
exact immersions for $\gamma_+$ with the leading derivatives of the
chosen loop.  We must show that their second forms preserve the later
curve conditions and that their preferred normals extend globally.
The crossing calculation below also specifies a turn threshold using
only the bounds that were fixed before the loop's large derivative.

\subsection{The exact second form}

Choose the loop of Lemma~\ref{lem:velocity-loop}, with the turn threshold
to be specified below, and let $\widetilde f_z$ be the resulting exact
immersions from Proposition~\ref{prop:primitive-realization}.  On the
disk their $yy$ second form has the representation
\[
 B_{\widetilde f_z}(\partial_y,\partial_y)=S n^*,\qquad
 |n^*|=1,\qquad S>0,
\]
for all sufficiently small $z$.  Indeed its limiting value is
$S_0n_0^*$ and $S_0$ has a positive minimum.  Project $m_0^*$ to the line
in the new normal plane perpendicular to $n^*$ and normalize; the
resulting $m^*$ is well defined and close to $m_0^*$.
Write
\[
 B_{\widetilde f_z}(\partial_x,\partial_y)=D^* n^*+N^* m^*,\qquad
 k=K_{\gamma_+}\det(\gamma_+).
\]
The determinant is in the $(x,y)$ coordinates.  The Gauss equation gives
\begin{equation}
 B_{yy}=(S,0),\qquad B_{xy}=(D^*,N^*),\qquad
 B_{xx}=\left(\frac{k+(D^*)^2+(N^*)^2}{S},\,L^*\right)
 \label{eq:exact-second-form}
\end{equation}
in the orthonormal frame $(n^*,m^*)$.  We have used
$\langle B_{xx},B_{yy}\rangle-|B_{xy}|^2=k$.
Thus the exact metric determines the first component of $B_{xx}$ from
the unscaled $yy$ and $xy$ data.  Only its second component, $L^*$,
requires the scaled $xx$ asymptotic from the realization theorem.
The derivative asymptotics of the exact realization imply, uniformly
on $\overline D$,
\begin{equation}
 \begin{gathered}
 S=S_0+o(1),\qquad n^*=n_0^*+o(1),\qquad
 D^*=D_0^*+o(1),\qquad N^*=N_0^*+o(1),\\
 zL^*=R_0\alpha_\theta+o(1).
 \end{gathered}
 \label{eq:second-form-limits}
\end{equation}
The final identity follows because
$V_\theta=R_0\alpha_\theta m_0^*$ is already normal to the limiting
pair $(X,Y)$.

For later use, write $Y_i=b_i\partial_x+c_i\partial_y$.  Formula
\eqref{eq:exact-second-form} gives the exact identities
\begin{equation}
 \begin{aligned}
 S\langle B(Y_i,Y_i),n^*\rangle
    &=(D^*b_i+Sc_i)^2+(k+(N^*)^2)b_i^2,\\
 \langle B(Y_i,Y_i),m^*\rangle
    &=L^*b_i^2+2N^*b_ic_i.
 \end{aligned}
 \label{eq:pure-direction-components}
\end{equation}
These two expressions separate the two mechanisms of the construction:
a large angular velocity gives a large $L^*$, whereas at a turn a large
transverse derivative gives a large $N^*$.

\subsection{Crossing thresholds without circular choices}

At a mutual later crossing inside $D$, both $b_i,b_l$ are nonzero.
Put
\[
 u_i=\partial_x+t_i\partial_y,\qquad
 u_l=\partial_x+t_l\partial_y,
 \qquad t_i=c_i/b_i,\quad t_l=c_l/b_l.
\]
The finitely many slopes are bounded.  From
\eqref{eq:exact-second-form}, the pure and mixed vectors are
\begin{equation}
 \begin{aligned}
 P_i:=B(u_i,u_i)
   &=\left(\frac{(N^*)^2+k+(D^*+St_i)^2}{S},\ L^*+2N^*t_i\right),\\
 P_l:=B(u_l,u_l)
   &=\left(\frac{(N^*)^2+k+(D^*+St_l)^2}{S},\ L^*+2N^*t_l\right),\\
 Z_{il}:=B(u_i,u_l)
   &=\left(\frac{(N^*)^2+k+(D^*+St_i)(D^*+St_l)}{S},\
             L^*+N^*(t_i+t_l)\right).
 \end{aligned}
 \label{eq:crossing-components}
\end{equation}
In particular
\begin{equation}
 Z_{il}-P_i=(t_l-t_i)(D^*+St_i,N^*).
 \label{eq:mixed-pure-difference}
\end{equation}

Choose bounds $0<s_-\le S\le s_+$, $|D^*|\le d_+$, and $|k|\le K_0$
on the entire closed disk and over all angular values, together with a
fixed bound on the finitely many crossing slopes.  They can be chosen before
the turn derivative: the limiting bounds for $S_0,1/S_0,D_0^*$ are
independent of that derivative, and after the loop is fixed,
\eqref{eq:second-form-limits} transfers slightly enlarged bounds to the
exact immersions for sufficiently small $z$.  There are constants
$C_1,C_2$ depending only on these bounds such that
\begin{equation}
 \begin{aligned}
 \langle P_j,n^*\rangle&\ge (N^*)^2/s_+-C_1,\\
 \left\langle Z_{il},\frac{P_j}{|P_j|}\right\rangle
 &\ge (N^*)^2/s_+-C_2(1+|N^*|),\qquad j=i,l,
 \end{aligned}
 \label{eq:uniform-crossing-bound}
\end{equation}
whenever the right side of the first line is positive.  To see the
second line, use
$\langle Z_{il},\widehat P_j\rangle
\ge |P_j|-|Z_{il}-P_j|$ and
\eqref{eq:mixed-pure-difference}.  In particular these estimates do not
depend on $L^*$.

Also
\[
 K_{\gamma_+}|u_i\wedge u_l|_{\gamma_+}^2
 =k(t_i-t_l)^2.
\]
We can therefore fix a finite threshold $T$ so large that $T^2>K_0+1$
and, when $|N^*|>T$, both pure vectors have positive dot with $n^*$ and
both mixed projections in \eqref{eq:uniform-crossing-bound} exceed
\[
 \sqrt{\max\{0,-k(t_i-t_l)^2\}}+1
\]
at every crossing.  Use a margin, for example $T+2$, as the turn
threshold in Lemma~\ref{lem:velocity-loop}, and then fix the loop.
No constant used to select this threshold depends on the resulting
large value of $(h_{\mathrm{turn}})_y$.  The sufficiently small choice
of $z$ may depend on this now fixed loop.

We next show positivity of both ordered crossing invariants for all
small $z$.  If a sequence of crossing phases tends to a turn,
\eqref{eq:second-form-limits} gives $|N^*|>T$, so the preceding estimate
applies.  If its limiting phase is not a turn, then
$|L^*|\to\infty$, while $N^*$ stays bounded for the fixed loop.
All three vectors in \eqref{eq:crossing-components} have the same
leading second component $L^*$; hence
\[
 \left\langle Z_{il},\widehat P_i\right\rangle\to+\infty,
 \qquad
 \left\langle Z_{il},\widehat P_l\right\rangle\to+\infty.
\]
These two cases and compactness of the phase circle give positivity at
the finitely many crossings uniformly for small $z$.
Returning to the original directions changes each ordered invariant
by the positive factor $b_i^2b_l^2$.  Indeed pure second forms scale by
squares, so their unit directions do not change even if either $b$ is
negative.  We have proved
\begin{equation}
 Q_{il}(\widetilde f_z)>0,\qquad
 Q_{li}(\widetilde f_z)>0
 \quad\text{at all later mutual crossings in }D.
 \label{eq:preserved-crossing-q}
\end{equation}

\subsection{Avoidance along the full curves}

We now prove that on every later curve inside the disk,
$B(Y_i,Y_i)$ is nonzero and its unit direction is not $-n^*$.
Suppose otherwise along a sequence $z\downarrow0$.
At a point of failure the second component in
\eqref{eq:pure-direction-components} is zero and the first is
nonpositive, so
\begin{equation}
 (D^*b_i+Sc_i)^2+(k+(N^*)^2)b_i^2\le0,
 \qquad L^*b_i^2+2N^*b_ic_i=0.
 \label{eq:opposite-ray-failure}
\end{equation}
Take a subsequence with limiting position and periodic phase.
The limiting $b_i$ cannot be zero: the vector field $Y_i$ is nonzero,
so then $c_i\ne0$ and the first inequality would give
$S_0^2c_i^2\le0$.  Multiplying the second equality by $z$ and using
\eqref{eq:second-form-limits} now forces
$\alpha_\theta=0$ at the limiting phase.

In the boundary collar, \eqref{eq:collar-inequality} rules out the
first inequality of \eqref{eq:opposite-ray-failure}, even without its
nonnegative $(N^*)^2b_i^2$ term.  On the remaining compact curve portions,
the limiting phase must be one of the two turns and
$k+(N_0^*)^2>0$ by the threshold just chosen.  This again contradicts
the first inequality.  Therefore, for all sufficiently small $z$,
\begin{equation}
 B(Y_i,Y_i)\ne0,
 \qquad \frac{B(Y_i,Y_i)}{|B(Y_i,Y_i)|}\ne-n^*
 \quad\text{on }C_i\cap\overline D.
 \label{eq:disk-ray-avoidance}
\end{equation}
This argument needs boundedness of $N^*$ only for the fixed loop while
$z$ tends to zero.  Such a bound may depend on the chosen turn size;
it is distinct from the derivative-independent bounds used to choose
that size.

\subsection{Adjusting the normal at crossings}

It remains to give both pure vectors positive dot product with a
preferred normal at each crossing.  If their dots with $n^*$ are
already positive, no change is needed.  At each crossing, choose small
neighborhoods of the two turning phases on which $|N_0^*|>T+1$.
The uniform limits and \eqref{eq:uniform-crossing-bound} make both
dot products positive on these neighborhoods for small $z$.
Consequently any crossing requiring adjustment has phase in their
compact complement, where $|\alpha_\theta|$ has a positive lower bound.
There $|L^*|\to\infty$ uniformly, and the two pure unit vectors are
arbitrarily close to the same sign of $m^*$.
For small $z$, denote this sign by $m$ and use the path
\begin{equation}
 n_t=\frac{(1-t)n^*+t\widehat P_i}
 {|(1-t)n^*+t\widehat P_i|},\qquad 0\le t\le1.
 \label{eq:normal-crossing-path}
\end{equation}
This path is defined and avoids the opposite of either pure unit
vector.  Its numerator has nonnegative $m$ component, strictly positive
for $t>0$, whereas the opposite pure vectors have negative $m$
component.  At $t=0$, use $n^*\perp m$.
At the endpoint both required dot products are positive because
$\widehat P_i$ and $\widehat P_l$ are close.

The vectors in this construction extend smoothly to a neighborhood
of the crossing.  Strictness, together with compactness of $[0,1]$,
gives a neighborhood on which the whole path remains defined and
avoids the opposite pure directions.  Replacing $t$ by a smooth
spatial cutoff, equal to one at the crossing, makes a smooth local
normal adjustment.  Since there are no triple intersections, its
support can avoid every other curve.  The finitely many crossing
neighborhoods can be chosen pairwise disjoint and disjoint from the
boundary collar.  These changes preserve
\eqref{eq:disk-ray-avoidance}; they do not affect
\eqref{eq:preserved-crossing-q}, whose quantities are independent of
the preferred normal.

\subsection{Extension through the exterior collar}

We finish with a global preferred normal.  The exact realization
converges to $F$ in $C^2$ on the exterior of $D$.  Thus the normalized
projection of the old normal $n$ to the new normal bundle defines a
smooth unit normal $n_{\mathrm{out}}$ there for small $z$.
On a thin exterior collar, the new $yy$ second form is nonzero and
its unit direction $n^*$ is close to the old direction
$\widehat A_0$.  By \eqref{eq:boundary-invariants},
$\widehat A_0$ and $n$ are not antipodal on $C_0$.
After narrowing the collar, the normalized convex interpolation
between $n^*$ and $n_{\mathrm{out}}$ is consequently defined.
Use it to match the unchanged $n^*$ near the disk boundary to
$n_{\mathrm{out}}$ farther outside.

We check the later curves during this interpolation.  At an old
crossing $C_i\cap C_0$, the Gauss interpretation
\eqref{eq:crossing-gauss-interpretation} gives
\begin{equation}
 \langle A_0(F),A_i(F)\rangle
 =|B_F(Y_0,Y_i)|^2+
 K_\gamma|Y_0\wedge Y_i|_\gamma^2
 \ge Q_{0i}(F)>0.
 \label{eq:old-pure-crossing-dot}
\end{equation}
Thus both old endpoint normals have positive dot product with the
later pure vector: one by this inequality, and the other by
\eqref{eq:boundary-invariants}.  Continuity and exterior $C^2$
convergence retain positivity near these crossings for the new
immersion, and convex interpolation preserves it.
A sufficiently thin collar meets later curves only in these crossing
neighborhoods.  It can avoid all later mutual crossings.
Away from the disk and collar, every old strict condition persists
by exterior $C^2$ convergence and compactness.  The resulting normal
$n_+$ is global and satisfies all of
\eqref{eq:boundary-invariants} for the later indices.

The second fundamental form is nonzero on the disk because $S>0$,
and it is nonzero on the exterior by $C^2$ convergence to $F$.
Choose $z$ small enough for these properties, all the finite crossing
conditions, and the prescribed exterior error $\varepsilon$.
The exact metric is $\gamma_+$ by
Proposition~\ref{prop:primitive-realization}.  This completes the
proof of Proposition~\ref{prop:primitive-addition}.

All frame orientations in this argument were chosen inside a
coordinate disk.  The global object carried from one step to the
next is only the unit section $n_+$ of the normal bundle.  The
construction therefore applies equally to orientable and
nonorientable surfaces.
